\chapter{Liczby zespolone}
\label{ma1:ch:complex-numbers}

W ciele $\R$ równanie $x^2+1=0$ nie ma rozwiązania, ponieważ kwadrat liczby rzeczywistej jest nieujemny. Skonstruujemy rozszerzenie ciała $\R$, w którym istnieje pierwiastek tego równania. Zamiast zakładać istnienie nowej liczby, zbudujemy ją jako parę liczb rzeczywistych. Następnie wyprowadzimy własności algebraiczne, geometryczne i podstawowe oszacowania potrzebne w analizie.

\section{Konstrukcja ciała liczb zespolonych}
\begin{definicja}
\label{ma1:def:complex-pairs}
Na zbiorze $\C=\R^2$ definiujemy dla $a,b,c,d\in\R$ działania
\[
 (a,b)+(c,d)=(a+c,b+d),\qquad
 (a,b)(c,d)=(ac-bd,ad+bc).
\]
Równość par oznacza równość obu współrzędnych.
\end{definicja}

\begin{twierdzenie}
\label{ma1:thm:complex-field}
Zbiór $\C=\R^2$ z działaniami z definicji \ref{ma1:def:complex-pairs} jest ciałem. Jego zerem jest $(0,0)$, jedynką $(1,0)$, elementem przeciwnym do $(a,b)$ jest $(-a,-b)$, a dla $(a,b)\ne(0,0)$ zachodzi
\[
 (a,b)^{-1}=\left(\frac{a}{a^2+b^2},-\frac{b}{a^2+b^2}\right).
\]
Odwzorowanie $\iota:\R\to\C$ dane przez $\iota(a)=(a,0)$ jest zanurzeniem ciał.
\end{twierdzenie}
\begin{proof}
Dodawanie jest łączne i przemienne, ponieważ odbywa się współrzędnie w $\R$. Dla dowolnego $(a,b)\in\C$ mamy $(a,b)+(0,0)=(a,b)$ i $(a,b)+(-a,-b)=(0,0)$. Dla dowolnych $(a,b),(c,d)\in\C$ zachodzi $(ac-bd,ad+bc)=(ca-db,cb+da)$, co dowodzi przemienności mnożenia, a $(1,0)(a,b)=(a,b)$.

Sprawdźmy łączność mnożenia. Dla $z=(a,b)$, $w=(c,d)$ i $u=(e,f)$ otrzymujemy
\[
\begin{aligned}
 (zw)u
 &=((ac-bd)e-(ad+bc)f,\ (ac-bd)f+(ad+bc)e)\\
 &=(ace-bde-adf-bcf,\ acf-bdf+ade+bce),\\
 z(wu)
 &=(a(ce-df)-b(cf+de),\ a(cf+de)+b(ce-df))\\
 &=(ace-adf-bcf-bde,\ acf+ade+bce-bdf).
\end{aligned}
\]
Obie współrzędne są równe. Rozdzielność wynika z
\[
\begin{aligned}
 z(w+u)&=(a(c+e)-b(d+f),a(d+f)+b(c+e))\\
 &=(ac-bd,ad+bc)+(ae-bf,af+be)=zw+zu.
\end{aligned}
\]
Rozdzielność z drugiej strony wynika z przemienności. Ponadto $(0,0)\ne(1,0)$.

Jeśli $(a,b)\ne(0,0)$, przynajmniej jeden z kwadratów $a^2,b^2$ jest dodatni, zatem $q=a^2+b^2>0$. Mnożenie daje $(a,b)(a/q,-b/q)=(1,0)$. Każdy element niezerowy ma więc odwrotność.

Wreszcie $\iota(a+c)=\iota(a)+\iota(c)$, $\iota(ac)=\iota(a)\iota(c)$ oraz $\iota(1)=(1,0)$. Jeśli $\iota(a)=\iota(c)$, to $a=c$, więc $\iota$ jest iniekcją.
\end{proof}

Utożsamiamy odtąd $\R$ z podciałem $\iota(\R)\subseteq\C$. Oznaczamy przez $\mathrm i$ parę $(0,1)$. Wówczas $\mathrm i^2=(-1,0)=-1$. Każda para $(a,b)$ jest równa $(a,0)+(b,0)(0,1)$, co zapisujemy jako $a+b\mathrm i$. Przedstawienie jest jednoznaczne: jeśli $a+b\mathrm i=c+d\mathrm i$, to z równości par wynika $a=c$ i $b=d$.

\begin{stwierdzenie}
\label{ma1:prop:complex-not-ordered}
Nie istnieje porządek na $\C$, względem którego $\C$ jest ciałem uporządkowanym.
\end{stwierdzenie}
\begin{proof}
W ciele uporządkowanym $1>0$, więc $-1<0$, natomiast kwadrat dowolnego elementu jest nieujemny. Równość $\mathrm i^2=-1$ prowadzi do sprzeczności.
\end{proof}

\begin{uwaga}
Konstrukcję można również opisać jako pierścień ilorazowy $\R[X]/(X^2+1)$. Istotnie, odwzorowanie $p\mapsto p(\mathrm i)$ zachowuje działania i jest surjektywne, ponieważ każda liczba zespolona jest postaci $a+b\mathrm i$. Dzieląc $p$ przez moniczny wielomian $X^2+1$, otrzymujemy $p=q(X)(X^2+1)+a+bX$. Stąd $p(\mathrm i)=a+b\mathrm i$, co jest zerem dokładnie wtedy, gdy $a=b=0$. Jądro jest więc ideałem $(X^2+1)$. Dwa wielomiany mają tę samą wartość w $\mathrm i$ dokładnie wtedy, gdy ich różnica należy do tego ideału. Przypisanie klasie wielomianu jego wartości w $\mathrm i$ jest zatem dobrze określone, iniektywne i surjektywne; zachowuje działania, więc jest izomorfizmem. Teorię pierścieni ilorazowych rozwijamy w Algebrze abstrakcyjnej.
\end{uwaga}

\section{Postać algebraiczna i sprzężenie}
\begin{definicja}
\label{ma1:def:complex-conjugation}
Dla $z=a+b\mathrm i\in\C$, gdzie $a,b\in\R$, definiujemy część rzeczywistą $\operatorname{Re}z=a$, część urojoną $\operatorname{Im}z=b$ oraz sprzężenie $\overline z=a-b\mathrm i$.
\end{definicja}

\begin{twierdzenie}
\label{ma1:thm:conjugation-properties}
Dla dowolnych $z,w\in\C$ zachodzą
\[
\begin{gathered}
 \overline{z+w}=\overline z+\overline w,\qquad
 \overline{zw}=\overline z\,\overline w,\qquad
 \overline{\overline z}=z,\\
 \operatorname{Re}z=\frac{z+\overline z}{2},\qquad
 \operatorname{Im}z=\frac{z-\overline z}{2\mathrm i},\qquad
 z\overline z=(\operatorname{Re}z)^2+(\operatorname{Im}z)^2.
\end{gathered}
\]
Ponadto $\overline z=z$ wtedy i tylko wtedy, gdy $z\in\R$, a dla $z\ne0$ mamy $z^{-1}=\overline z/(z\overline z)$ oraz $\overline{z^{-1}}=(\overline z)^{-1}$.
\end{twierdzenie}
\begin{proof}
Piszemy $z=a+b\mathrm i$, $w=c+d\mathrm i$. Z definicji
$\overline{z+w}=(a+c)-(b+d)\mathrm i=\overline z+\overline w$.
Ponadto
$\overline{zw}=(ac-bd)-(ad+bc)\mathrm i=(a-b\mathrm i)(c-d\mathrm i)$.
Dwukrotna zmiana znaku współrzędnej urojonej przywraca $z$.
Dodając i odejmując $z$ oraz $\overline z$, otrzymujemy odpowiednio $2a$ oraz $2b\mathrm i$. Mnożenie daje
$(a+b\mathrm i)(a-b\mathrm i)=a^2+b^2$.
Równość $\overline z=z$ jest równoważna $b=-b$, czyli $b=0$. Dla $z\ne0$ mamy $a^2+b^2>0$, a pomnożenie przez $\overline z/(a^2+b^2)$ daje jedynkę. Sprzężenie równości $zz^{-1}=1$ daje $\overline z\,\overline{z^{-1}}=1$, więc ostatnia własność wynika z jednoznaczności odwrotności.
\end{proof}
\begin{przyklad}
Niech $z=2+3\mathrm i$, $w=1-2\mathrm i$. Wtedy $z+w=3+\mathrm i$, $zw=8-\mathrm i$, a
\[
 \frac zw=\frac{(2+3\mathrm i)(1+2\mathrm i)}{1^2+2^2}
 =-\frac45+\frac75\mathrm i.
\]
Mnożenie ostatniej liczby przez $w$ daje $2+3\mathrm i$, co pozwala sprawdzić rachunek.
\end{przyklad}

\section{Moduł liczby zespolonej}
\begin{definicja}
\label{ma1:def:complex-modulus}
Dla $z=a+b\mathrm i\in\C$, gdzie $a,b\in\R$, modułem $z$ nazywamy liczbę
\[
 |z|=\sqrt{a^2+b^2}\in[0,+\infty).
\]
Pierwiastek jest dobrze określony przez twierdzenie \ref{ma1:thm:positive-root}.
\end{definicja}

\begin{twierdzenie}
\label{ma1:thm:complex-modulus}
Dla dowolnych $z,w\in\C$ zachodzą
\begin{enumerate}[label=(\roman*)]
\item $|z|\ge0$, przy czym $|z|=0$ wtedy i tylko wtedy, gdy $z=0$;
\item $|\overline z|=|z|$ oraz $|zw|=|z||w|$;
\item $|\operatorname{Re}z|\le |z|$, $|\operatorname{Im}z|\le |z|$;
\item $|z+w|\le |z|+|w|$;
\item $\big||z|-|w|\big|\le |z-w|$.
\end{enumerate}
Dla $w\ne0$ mamy również $|z/w|=|z|/|w|$.
\end{twierdzenie}
\begin{proof}
Piszemy $z=a+b\mathrm i$, $w=c+d\mathrm i$, gdzie $a,b,c,d\in\R$. Pierwsza własność wynika z nieujemności sumy kwadratów i jej zerowania się dokładnie dla $z=0$. Sprzężenie nie zmienia kwadratów współrzędnych. Korzystając z twierdzenia \ref{ma1:thm:conjugation-properties}, otrzymujemy
$|zw|^2=(zw)\overline{zw}=(z\overline z)(w\overline w)=|z|^2|w|^2$.
Obie liczby $|zw|$ i $|z||w|$ są nieujemne i mają ten sam kwadrat, więc są równe. Dla rzeczywistego $a$ moduł liczby zespolonej $a+0\mathrm i$ jest równy zwykłej wartości bezwzględnej $|a|$, gdyż obie liczby są nieujemne i mają kwadrat $a^2$. Ponieważ $a^2,b^2\le a^2+b^2$, a liczby $|a|$, $|b|$ i $|z|$ są nieujemne, ścisła monotoniczność kwadratu (lemat \ref{ma1:lem:power-bound}) daje $|a|\le |z|$ i $|b|\le |z|$.

Aby udowodnić nierówność trójkąta, obliczamy
\[
 |z+w|^2=|z|^2+2\operatorname{Re}(z\overline w)+|w|^2.
\]
Istotnie, po wymnożeniu $(z+w)(\overline z+\overline w)$ wyrazy mieszane sumują się do $z\overline w+\overline{z\overline w}$. Z (iii) i (ii) wynika
$\operatorname{Re}(z\overline w)\le |z\overline w|=|z||w|$.
Stąd $|z+w|^2\le (|z|+|w|)^2$, a obie strony przed podniesieniem do kwadratu są nieujemne. To daje (iv).
Stosując (iv) do $z=(z-w)+w$ oraz do $w=(w-z)+z$, otrzymujemy
$|z|-|w|\le |z-w|$ i $|w|-|z|\le |w-z|=|z-w|$, gdzie ostatnia równość wynika z $w-z=-(z-w)$ i $|-1|=1$. Razem dają (v).
Dla $w\ne0$ z $z=(z/w)w$ i (ii) wynika wzór na moduł ilorazu, ponieważ $|w|>0$.
\end{proof}

\begin{wniosek}
\label{ma1:cor:complex-triangle-equality}
Dla $z,w\in\C$ równość $|z+w|=|z|+|w|$ zachodzi wtedy i tylko wtedy, gdy przynajmniej jeden składnik jest zerem lub $z=tw$ dla pewnego $t>0$ w $\R$.
\end{wniosek}
\begin{proof}
Jeśli $z,w\ne0$, z dowodu nierówności trójkąta wynika, że równość zachodzi dokładnie wtedy, gdy $\operatorname{Re}(z\overline w)=|z||w|=|z\overline w|$. Dla $v=c+d\mathrm i\ne0$ równość $\operatorname{Re}v=|v|$ oznacza $c\ge0$ i $c^2=c^2+d^2$, czyli $d=0$ oraz $c>0$. Stąd $z\overline w$ jest dodatnią liczbą rzeczywistą. Ponieważ $w\overline w=|w|^2>0$, mamy $z=(z\overline w/|w|^2)w$. Odwrotnie, dla $t>0$ z multiplikatywności modułu otrzymujemy $|tw+w|=(t+1)|w|=|tw|+|w|$. Przypadki zerowe są bezpośrednie.
\end{proof}

\begin{przyklad}
Dla $z=3+4\mathrm i$ mamy $|z|=5$, natomiast $|z^{-1}|=1/5$. Dla $z=1$ i $w=\mathrm i$ zachodzi $|z+w|=\sqrt2<2=|z|+|w|$. Dla $z=1+\mathrm i$ i $w=2+2\mathrm i$ zachodzi równość.
\end{przyklad}

\section{Interpretacja geometryczna}
\label{ma1:sec:complex-geometry}
Liczbie $z=a+b\mathrm i$ odpowiada punkt $(a,b)$ płaszczyzny rzeczywistej, a także wektor od początku układu do tego punktu. Moduł jest długością tego wektora w sensie euklidesowym. Dodawanie liczb zespolonych odpowiada dodawaniu wektorów, a sprzężenie odbiciu względem osi rzeczywistej. Dla ustalonego $w$ odwzorowanie $z\mapsto z+w$ jest przesunięciem płaszczyzny o wektor odpowiadający $w$.

\begin{figure}[htbp]
\centering
\begin{tikzpicture}[scale=0.92,>=Stealth]
\draw[->] (-.6,0)--(4.4,0) node[right] {$\operatorname{Re}$};
\draw[->] (0,-2.6)--(0,2.8) node[above] {$\operatorname{Im}$};
\coordinate (Z) at (3,2);
\coordinate (B) at (3,-2);
\draw[dashed,gray] (Z)--(3,0)--(B);
\draw[dashed,gray] (Z)--(0,2) (B)--(0,-2);
\draw[->,blue!65!black,thick] (0,0)--(Z) node[above right] {$z=a+b\mathrm i$};
\draw[->,blue!65!black,thick] (0,0)--(B) node[below right] {$\overline z=a-b\mathrm i$};
\node[above left] at (1.5,1) {$|z|$};
\node[below left] at (0,0) {$0$};
\node[below right=2pt] at (3,0) {$a$};
\node[left] at (0,2) {$b$};
\node[left] at (0,-2) {$-b$};
\end{tikzpicture}
\caption{Sprzężenie jako odbicie względem osi rzeczywistej; przypadek $a,b>0$.}
\label{ma1:fig:complex-conjugation}
\end{figure}
\FloatBarrier

\begin{stwierdzenie}[Mnożenie jako odwzorowanie liniowe]
\label{ma1:prop:complex-multiplication-matrix}
Niech $z=a+b\mathrm i\in\C$ i $T_z:\C\to\C$, $T_z(w)=zw$. Przy utożsamieniu $\C$ z $\R^2$ odwzorowanie $T_z$ jest liniowe nad $\R$ i ma w bazie standardowej macierz
\[
 M_z=\begin{pmatrix}a&-b\\b&a\end{pmatrix}.
\]
Oznaczając przez $I_2$ macierz jednostkową stopnia $2$, otrzymujemy $M_z^\mathsf{T}M_z=|z|^2 I_2$, $\det M_z=|z|^2$ oraz $M_{zw}=M_zM_w$ dla $z,w\in\C$. Dla $z\ne0$ odwzorowanie $T_z$ jest bijekcją i skaluje wszystkie odległości euklidesowe przez $|z|$. Macierze ortogonalne stopnia $2$ o wyznaczniku $1$ są dokładnie macierzami $M_z$ dla $z\in S^1=\{v\in\C:|v|=1\}$.
\end{stwierdzenie}
\begin{proof}
Dla $w=c+d\mathrm i$ wzór na iloczyn daje współrzędne $(ac-bd,bc+ad)$, czyli iloczyn macierzy $M_z$ i kolumny $(c,d)^\mathsf{T}$. Liniowość wynika z liniowości mnożenia macierzowego. Rachunek daje
\[
 M_z^\mathsf{T}M_z=
 \begin{pmatrix}a&b\\-b&a\end{pmatrix}
 \begin{pmatrix}a&-b\\b&a\end{pmatrix}
 =\begin{pmatrix}a^2+b^2&0\\0&a^2+b^2\end{pmatrix}.
\]
Wyznacznik wynosi $a^2+b^2$. Równość $T_z(T_w(v))=(zw)v$ dla każdego $v\in\C$ wynika z łączności mnożenia, stąd $M_zM_w=M_{zw}$. Jeśli $z\ne0$, odwrotnością $T_z$ jest $T_{z^{-1}}$. Dla dowolnych $v,w\in\C$ mamy
$|T_z(v)-T_z(w)|=|z(v-w)|=|z||v-w|$.
 

Dla $|z|=1$ macierz $M_z$ jest ortogonalna o wyznaczniku $1$; takie odwzorowania płaszczyzny są obrotami wokół początku układu. Można to rozumieć bez funkcji trygonometrycznych: pierwsza kolumna $(a,b)$ jest jednostkowa, druga $(-b,a)$ jest do niej prostopadła, jednostkowa i zachowuje orientację bazy. Odwrotnie, jeżeli rzeczywista macierz ortogonalna stopnia $2$ ma wyznacznik $1$ i pierwszą kolumnę $(a,b)^\mathsf{T}$, to $a^2+b^2=1$. Wektor $(-b,a)^\mathsf{T}$ jest jednostkowy i prostopadły do pierwszej kolumny. Każdy wektor jednostkowy prostopadły do pierwszej kolumny jest równy temu wektorowi albo jego przeciwności (przestrzeń prostopadła ma wymiar $1$). Z warunku wyznacznika $1$ wynika wybór $(-b,a)^\mathsf{T}$, ponieważ wyznacznik macierzy o tych kolumnach to $a^2+b^2=1$, a zmiana znaku drugiej kolumny zmienia znak wyznacznika. Zatem macierz jest równa $M_{a+b\mathrm i}$. Przypisanie liczbowego kąta pozostawimy do czasu konstrukcji funkcji trygonometrycznych.
\end{proof}

\begin{figure}[htbp]
\centering
\begin{tikzpicture}[scale=1.35,>=Stealth]
\draw[->] (-1.6,0)--(1.85,0) node[right] {$\operatorname{Re}$};
\draw[->] (0,-1.5)--(0,1.85) node[above] {$\operatorname{Im}$};
\draw[gray] (0,0) circle (1);
\draw[->,gray,thick] (0,0)--(1,0) node[below right] {$1$};
\draw[->,gray,thick] (0,0)--(0,1) node[above left] {$\mathrm i$};
\draw[->,blue!65!black,thick] (0,0)--(.8,.6) node[right] {$z$};
\draw[->,blue!65!black,thick] (0,0)--(-.6,.8) node[above left] {$z\mathrm i$};
\draw[blue!65!black] (.12,.09)--(.03,.21)--(-.09,.12);
\node[below left] at (0,0) {$0$};
\end{tikzpicture}
\caption{Mnożenie przez $z=4/5+(3/5)\mathrm i$ zachowuje długości, prostopadłość i orientację.}
\label{ma1:fig:complex-rotation}
\end{figure}
\FloatBarrier

\section{Pierwiastki kwadratowe}
\begin{twierdzenie}
\label{ma1:thm:complex-square-root}
Każda liczba $w=u+v\mathrm i\in\C$, gdzie $u,v\in\R$, ma pierwiastek kwadratowy w $\C$. Dla $w=0$ jedynym pierwiastkiem jest $0$, a dla $w\ne0$ istnieją dokładnie dwa różne pierwiastki, będące elementami przeciwnymi. Jeśli $v\ne0$, jednym z nich jest $x+y\mathrm i$, gdzie
\[
 r=|w|,\qquad x=\sqrt{\frac{r+u}{2}},\qquad y=\frac{v}{2x}.
\]
Jeśli $v=0$ i $u>0$, pierwiastkami są $\pm\sqrt u$, natomiast jeśli $v=0$ i $u<0$, są nimi $\pm\mathrm i\sqrt{-u}$.
\end{twierdzenie}
\begin{proof}
Równanie $(x+y\mathrm i)^2=w$ jest równoważne układowi
$x^2-y^2=u$, $2xy=v$. Dla każdego rozwiązania $z=x+y\mathrm i$ z multiplikatywności modułu wynika dodatkowo $x^2+y^2=|z|^2=|z^2|=|w|=r$. Dodanie i odjęcie tych równań daje $x^2=(r+u)/2$ i $y^2=(r-u)/2$.
Dla $v\ne0$ zachodzi $r^2=u^2+v^2>u^2$, więc $r>|u|$. W szczególności $r+u>0$ i powyższa definicja $x$ daje $x>0$. Kładziemy $y=v/(2x)$. Ponieważ $v^2=r^2-u^2=(r-u)(r+u)$, otrzymujemy
\[
 y^2=\frac{v^2}{4x^2}
 =\frac{(r-u)(r+u)}{2(r+u)}=\frac{r-u}{2}.
\]
Stąd $x^2-y^2=u$ oraz $2xy=v$, czyli $(x+y\mathrm i)^2=w$. Dla $v=0$ i $u>0$ mamy $(\pm\sqrt u)^2=u$, a dla $v=0$ i $u<0$ zachodzi $(\pm\mathrm i\sqrt{-u})^2=(-1)(-u)=u$. Dla $w=0$ rozwiązaniem jest $0$.
Jeżeli $s^2=w$ i $t^2=w$, to $(s-t)(s+t)=0$. Ponieważ $\C$ jest ciałem, $s=t$ lub $s=-t$. Dla $w\ne0$ pierwiastek $s$ jest niezerowy, więc $s\ne-s$ (charakterystyka jest zerowa dzięki zanurzeniu $\R$). Dla $w=0$ z braku dzielników zera wynika $s=0$.
\end{proof}
\begin{przyklad}
Dla $w=3+4\mathrm i$ mamy $r=5$, $x=\sqrt{(5+3)/2}=2$ i $y=4/(2\cdot2)=1$. Zatem pierwiastkami są $2+\mathrm i$ i $-2-\mathrm i$. Dla $w=-9$ otrzymujemy $\pm3\mathrm i$. Dla $w=-3-4\mathrm i$ mamy $x=1$, $y=-2$, więc pierwiastkami są $1-2\mathrm i$ i $-1+2\mathrm i$.
\end{przyklad}

\begin{twierdzenie}[Równanie kwadratowe]
\label{ma1:thm:complex-quadratic}
Niech $A,B,C\in\C$, $A\ne0$. Niech $\Delta=B^2-4AC$ i niech $s\in\C$ spełnia $s^2=\Delta$. Zbiór rozwiązań równania $Az^2+Bz+C=0$ w $\C$ jest równy
\[
 \left\{\frac{-B+s}{2A},\frac{-B-s}{2A}\right\}.
\]
Rozwiązania są różne dokładnie wtedy, gdy $\Delta\ne0$.
\end{twierdzenie}
\begin{proof}
Istnienie $s$ wynika z twierdzenia \ref{ma1:thm:complex-square-root}. Dla każdego $z\in\C$ zachodzi tożsamość
\[
 (2Az+B)^2-\Delta=4A(Az^2+Bz+C).
\]
Ponieważ $4A\ne0$, pierwotne równanie jest równoważne równaniu $(2Az+B)^2=\Delta$. Twierdzenie \ref{ma1:thm:complex-square-root} daje $2Az+B\in\{s,-s\}$, a po rozwiązaniu względem $z$ podany zbiór. Różnica obu wyrażeń wynosi $s/A$, a więc jest niezerowa wtedy i tylko wtedy, gdy $s\ne0$, czyli $\Delta\ne0$.
\end{proof}
\begin{przyklad}
Równanie $z^2-2z+5=0$ ma wyróżnik $\Delta=-16$, stąd $s=4\mathrm i$ i rozwiązania $1\pm2\mathrm i$. Bezpośrednie podstawienie daje $(1\pm2\mathrm i)^2-2(1\pm2\mathrm i)+5=0$.
\end{przyklad}

\section{Okrąg jednostkowy i rozkład modułowy}
\begin{definicja}
\label{ma1:def:unit-circle}
Okręgiem jednostkowym na płaszczyźnie zespolonej nazywamy zbiór
$S^1=\{z\in\C:|z|=1\}$.
\end{definicja}
\begin{twierdzenie}
\label{ma1:thm:unit-circle-decomposition}
Zbiór $S^1$ jest podgrupą grupy multiplikatywnej $\C^\times=\C\setminus\{0\}$. Każde $z\in\C^\times$ ma dokładnie jedno przedstawienie $z=ru$, gdzie $r>0$ w $\R$ i $u\in S^1$. Mamy $r=|z|$, $u=z/|z|$, a odwzorowanie
\[
 \C^\times\to\R_{>0}\times S^1,\qquad z\longmapsto(|z|,z/|z|)
\]
jest izomorfizmem grup multiplikatywnych (działanie na iloczynie jest współrzędne).
\end{twierdzenie}
\begin{proof}
Jedynka należy do $S^1$. Jeśli $u,v\in S^1$, to $|uv|=|u||v|=1$, a $|u^{-1}|=1/|u|=1$. Stąd $S^1$ jest podgrupą. Ponadto $u^{-1}=\overline u$ dla $u\in S^1$, ponieważ $u\overline u=|u|^2=1$.
Dla $z\ne0$ mamy $|z|>0$ i $\left|z/|z|\right|=|z|/|z|=1$, co daje istnienie przedstawienia. Jeśli $z=ru$ przy podanych założeniach, to $|z|=|r||u|=r$ i $u=z/|z|$, co daje jednoznaczność.
Multiplikatywność modułu oraz przemienność mnożenia dają dla $z,w\ne0$
\[
 (|zw|,zw/|zw|)
 =(|z||w|,(z/|z|)(w/|w|)).
\]
Odwzorowanie zachowuje więc działanie. Jest bijekcją, ponieważ odwzorowanie $(r,u)\mapsto ru$ jest jego odwrotnością.
\end{proof}
\begin{uwaga}
Rozkład $z=ru$ nie wymaga jeszcze funkcji wykładniczej ani trygonometrycznych. Po ich skonstruowaniu wyrazimy $u$ za pomocą argumentu i otrzymamy postać trygonometryczną, wzór Eulera oraz wzory na pierwiastki wyższych stopni. Nie korzystamy z tych wyników przed ich udowodnieniem. W następnym rozdziale przejdziemy do ciągów w przestrzeniach metrycznych; moduł liczby zespolonej będzie jednym z podstawowych przykładów konstrukcji odległości.
\end{uwaga}
